\chapter{Varifolds and first variation}\label{sec:varifolds}
A varifold records location and an unoriented tangent plane. This formulation places smooth and nonsmooth surfaces in a common space of Radon measures. We use \cite{Sim18}, Chapter 8, Sections 1--3.

\begin{definition}[Varifold and weight]\label{def:varifold}
Let $1\leq m<n$. The Grassmannian $G(n,m)$ is the space of unoriented $m$-dimensional linear subspaces of $\mathbb{R}^n$, with the topology induced by the orthogonal projection matrices. An \emph{$m$-varifold in $U$} is a nonnegative Radon measure on $G_m(U)=U\times G(n,m)$. Its weight is the Radon measure on $U$ defined by
\[
 \|V\|(A)=V(A\times G(n,m))\qquad(A\subset U\text{ Borel}).
\]
Varifold convergence means weak convergence of these measures against $C_c(G_m(U))$. See \cite{Sim18}, Chapter 8, Section 1, pp.~235--236.
\end{definition}

\begin{definition}[Rectifiable and integral varifolds]\label{def:rectifiable-varifold}
Let $M\subset U$ be $\mathcal{H}^m$-measurable and countably $m$-rectifiable, with locally finite $\mathcal{H}^m$ measure in $U$. Let $\theta:M\to(0,\infty)$ be measurable and satisfy $\int_{M\cap K}\theta\,d\mathcal{H}^m<\infty$ for every compact $K\subset U$. The associated rectifiable varifold $V(M,\theta)$ is determined by
\[
 \int\varphi(x,S)\,dV(M,\theta)(x,S)
 =\int_M\varphi(x,T_xM)\theta(x)\,d\mathcal{H}^m(x).
\]
It is integral if $\theta$ is integer-valued almost everywhere. Its weight is $\theta\mathcal{H}^m\llcorner M$. The tangent planes exist almost everywhere by Theorem~\ref{thm:rectifiable-tangents}; changing the integrand on the exceptional null set has no effect. See \cite{Sim18}, Chapter 4, Section 1, and Chapter 8, p.~236.
\end{definition}

\begin{definition}[First variation, stationarity, and mean curvature]\label{def:first-variation}
For an $m$-varifold $V$ and $X\in C_c^1(U;\mathbb{R}^n)$, define
\[
 \delta V(X)=\int_{G_m(U)}\operatorname{div}_S X(x)\,dV(x,S),\qquad
 \operatorname{div}_S X=\sum_{i=1}^m\langle D_{e_i}X,e_i\rangle,
\]
where $(e_i)$ is any orthonormal basis of $S$. The trace is independent of that basis. We call $V$ stationary in $U$ if $\delta V(X)=0$ for every such $X$.

It has locally bounded first variation if for each $W\Subset U$ there is $C_W<\infty$ such that $|\delta V(X)|\leq C_W\sup|X|$ whenever $X$ is supported in $W$. In the case when the first variation is represented entirely by a vector field $H\in L^1_{\mathrm{loc}}(\|V\|;\mathbb{R}^n)$,
\[
 \delta V(X)=-\int_U H\cdot X\,d\|V\|,
\]
we call $H$ the generalized mean curvature. Locally bounded first variation alone may also have a singular boundary term; the last identity assumes that term is absent. See \cite{Sim18}, Chapter 8, equations (2.3)--(2.10), pp.~239--240.
\end{definition}

\begin{proposition}[An affine plane is stationary]\label{prop:plane-stationary}
Let $P\subset\mathbb{R}^n$ be an affine $m$-plane and let $\theta>0$ be constant. Then $V(P,\theta)$ is stationary in $\mathbb{R}^n$.
\end{proposition}
\begin{proof}
Identify $P$ isometrically with $\mathbb{R}^m$. The restriction of a compactly supported vector field has compact support on $P$, and its tangential divergence is the Euclidean divergence of its tangential component. The integral of each coordinate derivative is zero by the fundamental theorem of calculus and Fubini's theorem. Hence $\delta V(X)=\theta\int_P\operatorname{div}_P X\,d\mathcal{H}^m=0$.
\end{proof}

\begin{theorem}[Stationary monotonicity]\label{thm:stationary-monotonicity}
Let $V$ be a stationary $m$-varifold in $U$, let $a\in U$, and suppose $\overline{B_R(a)}\subset U$. Then
\[
 r\longmapsto\frac{\|V\|(B_r(a))}{\omega_mr^m}
 \quad\text{is nondecreasing for }0<r<R.
\]
Consequently $\Theta^m(\|V\|,a)$ exists and is finite and nonnegative. No rectifiability hypothesis is needed for this conclusion, and the density may be zero.

\emph{Cited monotonicity theorem; proof to be added.} See \cite{Sim18}, Chapter 8, Section 3, equations (3.5)--(3.6), p.~241, with $\Lambda=0$. Finiteness of the limiting density follows by bounding all sufficiently small ratios by the ratio at any fixed radius less than $R$ and using local finiteness of $\|V\|$.
\end{theorem}
